\chapter{CONCLUSIONES}

De acuerdo con los objetivos planteados y la contrastación empírica de las hipótesis de la investigación, se presentan las siguientes conclusiones:

\noindent \textbf{PRIMERA:} Se logró desarrollar y evaluar experimentalmente un modelo de clasificación híbrido basado en la arquitectura de ensamble por Stacking de dos niveles sobre el dataset CICIoT2023, integrando Random Forest, XGBoost GPU y LightGBM con cuantificación de incertidumbre por Entropía de Shannon y sobremuestreo adaptativo Borderline-SMOTE1. El modelo propuesto alcanzó un rendimiento global de $70.18\%$ en la métrica F1-Macro y una exactitud (Accuracy) del $84.85\%$ sobre la partición de prueba aislada de 377,383 observaciones, superando a todos los clasificadores de control individuales (XGBoost 69.60\%, LightGBM 69.53\%, Random Forest 67.91\%, CART 66.51\% y Regresión Logística 42.20\%). Esta superioridad predictiva fue confirmada estadísticamente mediante la Prueba de Rangos con Signo de Wilcoxon ($p = 0.003744 < 0.05$), rechazando la hipótesis nula y aceptando formalmente la Hipótesis General de la investigación.

\vspace{0.8em}
\noindent \textbf{SEGUNDA:} Se procesó el conjunto de datos de red IoT de forma exitosa mediante la implementación del escalado robusto adaptativo (\textit{RobustScaler}) y el sobremuestreo sintético en la frontera (\textit{Borderline-SMOTE1}) sobre el subconjunto de entrenamiento, logrando neutralizar el sesgo algorítmico provocado por el desbalance extremo de clases ($\mathrm{IR} = 41.75:1$). Esta etapa de preprocesamiento permitió elevar la sensibilidad (\textit{Recall}) y el F1-Score en los ataques minoritarios infrecuentes a 41.40\% en la clase \textit{Brute Force} y a 33.86\% en la clase \textit{Web-based Attacks}, mitigando la paradoja de la exactitud y superando el colapso predictivo exhibido por los modelos tradicionales sin rebalanceo, aceptando la Hipótesis Específica 1.

\vspace{0.8em}
\noindent \textbf{TERCERA:} Se diseñó y entrenó la arquitectura de ensamble por Stacking integrando modelos base heterogéneos de Nivel 0 enriquecidos con la cuantificación de incertidumbre entrópica de Shannon ($H_{\mathrm{RF}}, H_{\mathrm{XGB}}, H_{\mathrm{LGB}}$) y sintonizando un meta-modelo de Nivel 1 basado en LightGBM mediante \texttt{GridSearchCV}. La inclusión de los descriptores entrópicos en el vector de meta-características $Z \in \mathbb{R}^{27}$ proporcionó al meta-modelo un indicador explícito sobre la ambigüedad probabilística ($H \to 3.0\text{ bits}$), permitiendo corregir eficazmente las clasificaciones erróneas en la franja limítrofe entre tráficos benignos y vectores maliciosos, alcanzando un F1-Macro del 70.18\% y aceptando la Hipótesis Específica 2.

\vspace{0.8em}
\noindent \textbf{CUARTA:} Se evaluó el rendimiento computacional del modelo de Stacking Híbrido demostrando una latencia de inferencia de apenas $8.65\,\mu s$ por muestra (equivalente a una tasa de procesamiento de 115,613 paquetes por segundo) y un consumo de memoria principal de $257.89\text{ MB}$, manteniéndose estrictamente por debajo del umbral operativo de $300\text{ MB}$ de memoria RAM exigido en dispositivos de borde. Esto demuestra que la arquitectura propuesta ofrece plena viabilidad computacional para su ejecución en tiempo real en pasarelas IoT frente a las alternativas pesadas de Deep Learning, convalidando la Hipótesis Específica 3.